\section{Security Considerations}

\subsection{Threat model}

Each layer's threat model is its own; the ZAP Stack does not
introduce a new aggregate threat model. The composition is
load-bearing only insofar as no layer's invariants are violated by
the layer's interaction with its neighbors.

\paragraph{Layer~0 (frame).} The wire frame is a magic byte plus a
length-prefixed object. A malformed frame is rejected before parsing.
Tests for the parser are in \texttt{luxfi/zap}; rejecting bad-magic
inputs is the first invariant
(\texttt{/tmp/zap-vs-codec/claim1\_reject\_bad\_magic\_test.go}).

\paragraph{Layer~3 (identity).} \texttt{sha256.Sum256(buf)} is
collision-resistant under the standard SHA-2 assumptions. The
identity is over the canonical byte stream; there is no second
canonical form.

\paragraph{Layer~4 (transcript).} TupleHash256 with NIST SP 800-185
framing is collision-resistant under cSHAKE256. The framing is
unambiguous: typed field separators prevent length-extension and
type-confusion attacks on the digest.

\paragraph{Layer~5 (signature).} \texttt{ecdsa(sha256(unsigned\_bytes))}
is a standard ECDSA construction. The signed window is precisely the
unsigned-bytes portion of the envelope; the credentials window is
appended after signing and does not feed back into the signature
domain.

\paragraph{Layer~9 (consensus cert).} The QuasarCert (LP-182) composes
legs as opaque byte strings; each leg is verified independently. The
security of the composition is the maximum of the leg securities --
a relying party at PQ-strict (BLS + Pulsar + Corona) survives any
single-leg break.

\subsection{What does NOT change with ZAP-native}

The cryptographic primitives are unchanged. ZAP-native changes how
bytes flow between layers; it does not change the underlying ECDSA,
BLS, ML-DSA, Corona, or SLH-DSA constructions. Any security claim
proven against the legacy stack is preserved by the ZAP Stack iff
that claim holds against the same primitive consumed via offset-read
instead of via reflective marshal -- which is to say, always, because
the bytes are identical.

\subsection{Side channels}

The offset-read pattern uses constant-position field access; there
is no data-dependent branching in the parser. The signature verifier
operates on the byte window; it does not allocate fresh structs that
the GC could observe. The reduction in allocations (Table~\ref{tab:perops},
``2 allocs $\to$ 1 alloc'') closes a side channel that the legacy
\texttt{codec.Manager} held open: the codec registry lookup was
data-dependent in the type tag.

\section{Backwards Compatibility}

None. The activation marker is at the genesis of the new final Lux
network: unix \texttt{1766708400}. The pre-Quasar Edition Lux network
(2020--2025) is a separate network operating a different stack. There
is no replay, no flag day, no migration path.
